% !TeX program = lualatex
% Material complementar — documento do aluno.
% Questões: as 19 marcadas "C" em fontes/enem-25-selecionadas.md e a A5 (2014_CH_038), que saiu da aula;
% recortes do PDF oficial em fontes/provas-recortes/.
% Compila com pdfLaTeX ou LuaLaTeX. A opção [final] faz qualquer pendência virar erro.
\documentclass[11pt]{article}
\usepackage[final]{aula-geo}
../estilo/metadados.tex
\definir\tipodocumento{Material complementar}
\definir\versao{v3}
\graphicspath{{../../../fontes/provas-recortes/}}

\begin{document}
\blocotitulo

\section*{Como usar este material}
Este material continua a aula \emph{Onde cabe um data center?}. Na aula, o mapa foi lido duas vezes: primeiro
como o planejador o lê (onde há energia, transmissão, acesso e terreno) e depois do avesso (quem e o que está
nesse mapa e não aparece na conta). As questões abaixo treinam as mesmas ferramentas em outros assuntos.

\begin{enumerate}
\item Comece refazendo as quatro questões resolvidas em aula, sem olhar as anotações. Depois, escolha um grupo
      por dia. Cada grupo começa com uma frase do que retomar da aula.
\item Resolva sem consultar. Marque as que travaram, porque o erro diz mais que o acerto.
\item Só depois confira no gabarito, no fim. Para cada erro, escreva em uma linha \emph{por que} a alternativa certa
      é certa; não basta copiar a letra.
\item Depois dos grupos, faça a \emph{Proposta de redação}: ela reaproveita o que apareceu na aula e nas questões.
\end{enumerate}

\begin{destaque}[title={Antes de começar}]
Todas as questões são reais, do caderno azul do 1º dia, de 2013 a 2016. Nenhum assunto aqui exige saber algo sobre
data centers: eles são só o pretexto da aula. O que o ENEM cobra é ler texto, mapa, gráfico e tabela.
\end{destaque}

\clearpage
\section{Questões do ENEM}
Questões reais do ENEM (2013--2016), reproduzidas da prova oficial do INEP (caderno azul, 1º dia).
A numeração deste material vem antes de cada questão; o número dentro da imagem é o da prova.
O gabarito está no fim. As questões 1 a 4 foram resolvidas em aula; da 5 em diante, são novas.

\Needspace{0.6\textheight}
\subsection*{Resolvidas em aula}
\emph{Retome:} refaça sozinho e diga por que as outras alternativas estão erradas.
\begin{multicols}{2}\raggedcolumns

\questaoenem{2013_CH_043}{2013}{43}{C}
\questaoenem{2014_CH_009}{2014}{9}{E}
\questaoenem{2015_CH_015}{2015}{15}{D}
\questaoenem{2013_CH_037}{2013}{37}{C}
\end{multicols}

\subsection*{Energia, técnica e informação}
\emph{Retome:} matriz energética $\times$ matriz elétrica; objeto técnico; o que a técnica distribui e o que ela concentra.
\begin{multicols}{2}\raggedcolumns

\questaoenem{2014_CH_034}{2014}{34}{C}
\questaoenem{2015_CH_029}{2015}{29}{E}
\questaoenem{2016_CH_040}{2016}{40}{C}
\questaoenem{2015_CH_021}{2015}{21}{E}
\questaoenem{2015_CH_038}{2015}{38}{D}
\end{multicols}

\Needspace{0.6\textheight}
\subsection*{Cidade, rede e mapa}
\emph{Retome:} camada de mapa; fatores locacionais; por que serviços se concentram em polos e como a rede urbana se reorganiza.
\begin{multicols}{2}\raggedcolumns

\questaoenem{2013_CH_034}{2013}{34}{C}
\questaoenem{2016_CH_004}{2016}{4}{C}
\questaoenem{2013_CH_001}{2013}{1}{E}
\questaoenem{2014_CH_026}{2014}{26}{D}
\questaoenem{2015_CH_012}{2015}{12}{E}
\end{multicols}

\Needspace{0.6\textheight}
\subsection*{Água}
\emph{Retome:} água como recurso que muda de uso, de disponibilidade e de quem decide; leitura de tabela e gráfico.
\begin{multicols}{2}\raggedcolumns

\questaoenem{2014_CH_003}{2014}{3}{B}
\questaoenem{2014_CH_038}{2014}{38}{E}
\questaoenem{2015_CH_032}{2015}{32}{E}
\questaoenem{2016_CH_036}{2016}{36}{C}
\end{multicols}

\Needspace{0.6\textheight}
\subsection*{Clima, relevo e biomas}
\emph{Retome:} a ficha de lugar (domínio, relevo, clima, hidrografia, impacto) e a leitura de climograma e de mapa de biomas.
\begin{multicols}{2}\raggedcolumns

\questaoenem{2014_CH_035}{2014}{35}{E}
\questaoenem{2015_CH_030}{2015}{30}{B}
\questaoenem{2015_CH_027}{2015}{27}{D}
\questaoenem{2014_CH_036}{2014}{36}{D}
\questaoenem{2015_CH_001}{2015}{1}{A}
\questaoenem{2016_CH_030}{2016}{30}{E}
\end{multicols}

\clearpage
\section{Proposta de redação}
Treino no formato do ENEM. Os textos motivadores foram elaborados para este material com base nas fontes indicadas.

\subsection*{Textos motivadores}

\textbf{Texto I}\par\nopagebreak
\begin{quote}
\enquote{A nuvem, que é a espinha dorsal da indústria da inteligência artificial, é formada por pedras, salmouras de
lítio e petróleo bruto.}
\end{quote}
\fonte{CRAWFORD, K. \emph{Atlas da I.\,A.} São Paulo: Edições Sesc São Paulo, 2025, p.~43.}

\medskip
\textbf{Texto II}\par\nopagebreak
Em 2024, as fontes renováveis responderam por 88,2\% da matriz elétrica brasileira. A eletricidade abundante e de
origem renovável é apresentada pelo setor como vantagem do país para atrair data centers, os prédios onde ficam os
computadores que armazenam dados e fazem funcionar a inteligência artificial.
\fonte{elaborado com base em EPE. \emph{Balanço Energético Nacional 2025}: relatório síntese. Rio de Janeiro: EPE, 2025, p.~38.}

\medskip
\textbf{Texto III}\par\nopagebreak
Em setembro de 2024, meses depois da enchente que atingiu o Rio Grande do Sul, a empresa Scala anunciou uma
\enquote{cidade de data centers} em Eldorado do Sul, na região metropolitana de Porto Alegre, com investimento inicial
de R\$~3 bilhões, segundo a empresa. O data center que a mesma empresa mantinha em Porto Alegre ficou 68 dias fora
do ar durante a enchente. O governo gaúcho afirmou não ter concedido incentivo fiscal especial ao projeto.
\fonte{elaborado com base em SETTI, R. Depois da tragédia, a inteligência artificial. \emph{O Globo}, 11 set. 2024.}

\medskip
\textbf{Texto IV}\par\nopagebreak
Cerca de 8 em cada 10 moradores de Eldorado do Sul foram diretamente atingidos pela enchente de 2024, a maior
proporção do Rio Grande do Sul. Perto do terreno escolhido para o empreendimento ficava a aldeia Tekoa Pekuruty,
do povo Mbyá-Guarani, destruída depois da enchente. Seus moradores esperam realocação e dizem não ter sido
consultados sobre o projeto. A Convenção 169 da Organização Internacional do Trabalho prevê a consulta prévia,
livre e informada aos povos indígenas sobre medidas que os afetem.
\fonte{elaborado com base em IPEA. \emph{Uma estimativa da população atingida pelas enchentes do Rio Grande do Sul em 2024}.
Brasília: Ipea, 2024; e MARTINS, L. \emph{Intercept Brasil}, 23 jun. 2025.}

\begin{redacao}[Proposta]
A partir da leitura dos textos motivadores e com base nos conhecimentos construídos ao longo de sua formação,
redija um texto dissertativo-argumentativo em modalidade escrita formal da língua portuguesa sobre o tema
\textbf{\enquote{Desafios da expansão dos data centers no território brasileiro}}, apresentando proposta de
intervenção que respeite os direitos humanos. Selecione, organize e relacione, de forma coerente e coesa,
argumentos e fatos para a defesa de seu ponto de vista.
\end{redacao}

\begin{destaque}[title={Instruções}]
O texto definitivo tem até 30 linhas; use a folha de redação no fim deste material. Recebe nota zero o texto com
até 7 linhas, o que foge ao tema e o que não é dissertativo-argumentativo. Linhas copiadas dos textos motivadores
não contam.
\end{destaque}

\Needspace{18\baselineskip}
\begin{redacao}[Antes de escrever: monte a proposta de intervenção]
Escreva, em uma frase cada, os cinco elementos exigidos pela Competência 5 do ENEM:
\begin{enumerate}
\item \textbf{Agente:} quem age (ministério, órgão ambiental, prefeitura, empresa, escola).
\item \textbf{Ação:} o que essa pessoa ou instituição faz.
\item \textbf{Meio:} por qual instrumento ou modo (lei, fiscalização, campanha, audiência pública).
\item \textbf{Finalidade:} para quê, com verbo no infinitivo.
\item \textbf{Detalhamento:} um dado a mais sobre o agente, a ação, o meio ou a finalidade.
\end{enumerate}
Depois, junte as cinco frases em um parágrafo e troque com um colega: ele deve achar cada elemento sem perguntar.
\end{redacao}

\subsection*{Repertório da aula}
\begin{itemize}
\item A \enquote{nuvem} é material e tem endereço: energia, água, terra e rede (Texto I).
\item Objeto técnico e guerra dos lugares, conceitos de Milton Santos (\emph{A natureza do espaço}), aplicados ao caso:
      os lugares competem para atrair empreendimentos.
\item As duas leituras do mapa: a do planejador (energia, transmissão, acesso, terreno) e a de quem vive no lugar.
\item Consulta prévia, livre e informada (Convenção 169 da OIT).
\end{itemize}

\subsection*{Depois de escrever: confira pelas cinco competências}
Cada competência vale de 0 a 200 pontos. Releia o seu texto e responda:
\begin{enumerate}[label=\textbf{C\arabic*}]
\item Escrevi na norma-padrão, sem erros de grafia, concordância, regência e pontuação?
\item Tratei do tema inteiro (desafios, expansão, território brasileiro)? O texto é dissertativo-argumentativo e
      usa repertório de outras áreas a serviço do argumento?
\item Tenho uma tese clara e argumentos que a defendem, em ordem, sem repetir ideias?
\item Os parágrafos e as frases se ligam com conectivos variados e bem empregados?
\item A proposta de intervenção tem os cinco elementos, respeita os direitos humanos e responde ao problema que eu discuti?
\end{enumerate}

\clearpage
\subsection*{Folha de redação}
\newcount\linharedacao \linharedacao=0
\vspace{0.6em}
{\setlength{\parskip}{0pt}%
\loop\advance\linharedacao by 1
  \noindent\makebox[1.6em][r]{\sffamily\scriptsize\color{fgsoft}\the\linharedacao}\hspace{0.4em}%
  {\color{regua}\rule{\dimexpr\linewidth-2em\relax}{0.4pt}}\par\vspace{0.22cm}%
\ifnum\linharedacao<30\repeat}

\clearpage
\section*{Gabarito}
\imprimirgabarito

\end{document}
